\chapter{Introduction}
\label{ch:intro}

% GUIDANCE (chapter target about 1,500 words):
% Open with the festival scene, then move to the measurement problem in two paragraphs. Do not open
% with "In recent years". Cite the T2I lineage once (\cite{ramesh2021zero, ramesh2022hierarchical,
% betker2023improving}) and the evaluation literature once (\cite{hessel2021clipscore, hu2023tifa,
% ghosh2023geneval, lin2024evaluating}) to establish that existing metrics score objects, not affect.

\section{Context}
\label{sec:intro-context}
% GUIDANCE: The corpus: \nTotal\ prompt--drawing records collected at Festival of the Mind 2024 in
% Sheffield; visitors were invited to describe how they were feeling, the interface accepted free
% text, and a model drew each entry as a black-and-white line drawing. Records are not linked to
% people: \nRepeats\ repeat an earlier prompt, so the number of contributors is unknown and no greater
% than \nDistinct. State that a substantial minority (\nNonEmotional\ records by the prompt reader)
% are object or scene requests ("pigs") rather than self-report, and that this is a property of the
% collection, not a defect of the analysis. Keep it factual; no numbers beyond the macros.

\section{The question}
\label{sec:intro-question}
% GUIDANCE: One question, stated once: when a text-to-image model draws what a person wrote about
% their feeling, does the emotion recoverable from the drawing agree with the emotion recoverable
% from the words? Then the supervisor's reframing, which is the intellectual centre of the project:
% SIMILARITY, NOT MATCHING. Emotion labels overlap and blend (\cite{cowen2017self, russell1980circumplex,
% posner2005circumplex}), so "happy" versus "excited" is a near miss, not an error; the right
% instrument is a ladder of measures of increasing strictness, each read against what random pairing
% would give. Introduce the four rungs by name only (valence--arousal similarity, valence sign,
% family, top-3 overlap) and point to Chapter~\ref{ch:method}.
% Also say here what the project does NOT claim, in one sentence: nothing in it is a statement about
% what a person would see in a drawing — agreement is between two machine readers, and the human
% study is future work.

\section{Aims and objectives}
\label{sec:intro-aims}
% GUIDANCE: One aim, then objectives as a numbered list (the marker maps the Conclusions back to this):
%   1. Build a pipeline that reads an emotion from a monochrome drawing WITHOUT asking a vision model
%      for one: describe with factual questions, then extract from the description.
%   2. Define a similarity ladder and a random-pairing null so that every value has a chance level.
%   3. Measure agreement on the whole corpus and on the splits the supervisor asked for
%      (emotional vs neutral-labelled prompts; face vs scene drawings).
%   4. Test the face/scene difference formally and check robustness to repeated prompts.
%   5. Compare the describe-then-extract reader with a direct CLIP reader like for like.
%   6. Use the repeated prompts to measure the stability of the whole prompt-to-reading chain.
%   7. Audit the instrument: the control set, the describer, an error taxonomy, an ablation of the
%      description template, and a single-coder reliability check.
%   8. Release a reproducible toolkit with pinned numeric tests and a live demonstrator.

\section{Contributions}
\label{sec:intro-contributions}
% GUIDANCE: Honest novelty accounting (from RESEARCH_GAP.md): what is new is the OBJECT (unmediated,
% mixed-affect public self-report rendered as colour-free line art), the QUESTION (end-to-end affect
% agreement on non-templated input, cf. \cite{lomas2024improved} who use a fixed typology), the
% NULL-REFERENCED similarity ladder, and the CONSISTENCY experiment that the repeated prompts make
% possible. What is incremental is the toolkit (existing models recombined). Say so.
% List the deliverables: the analysis code, \nAnalysed-record result tables, the demonstrator, the
% test suite, and the audits in Chapter~\ref{ch:results}.

\section{Overview of the report}
\label{sec:intro-overview}
% GUIDANCE: One sentence per chapter. Chapter~\ref{ch:lit} reviews the literature in eight themes and
% states the gap. Chapter~\ref{ch:method} gives the corpus, pre-processing, pipeline, measures, null,
% inference and audits. Chapter~\ref{ch:results} reports the results, tables first. Chapter~\ref{ch:discussion}
% interprets them, says what is not claimed, and relates them to prior work. Chapter~\ref{ch:conclusion}
% concludes. Appendices give reproducibility and the instruments.
